% SEGMENT OLP-0065-S01
% Part: first-order-logic
% Chapter: proof-systems
% Section: seqeunt-calculus

\documentclass[../../../include/open-logic-section]{subfiles}

\begin{document}

\iftag{FOL}
      {\olfileid{fol}{prf}{seq}}
      {\olfileid{pl}{prf}{seq}}

\olsection{தொடரணிக் கணியம்}

பல !!{derivation} முறைமைகள் !!{sentence}s இன் ஒழுங்கமைப்புகளைக் கொண்டு
செயல்படுகின்றன; தொடரணிக் கணியம் \emph{தொடரணிகளைக்} கொண்டு செயல்படுகிறது.
ஒரு தொடரணி பின்வரும் வடிவிலான கோவை:
\[
!A_1, \dots, !A_m \Sequent !B_1, \dots, !B_m,
\]
அதாவது, தொடரணிக் குறி~$\Sequent$ ஆல் பிரிக்கப்பட்ட !!{sentence}s இன்
இரண்டு தொடர்களைக் கொண்ட சோடி. இவ்விரு தொடர்களில் எதுவும் வெறுமையாக
இருக்கலாம். தொடரணிக் கணியத்தில் !!^a{derivation} என்பது தொடரணிகளின் மரம்.
அதன் உச்சியிலுள்ள தொடரணிகள் சிறப்பு வடிவில் இருக்கும் (அவை ``தொடக்கத்
தொடரணிகள்'' அல்லது ``அடிகோள்கள்'' எனப்படுகின்றன); மற்ற ஒவ்வொரு தொடரணியும்
தனக்கு உடனடியாக மேலுள்ள தொடரணிகளிலிருந்து ஓர் அனுமான விதியால் பின்வரும்.
அனுமான விதிகள், தொடரணிகளிலுள்ள !!{sentence}s ஐ இடப்புறத்திலோ
வலப்புறத்திலோ சேர்த்து, நீக்கி அல்லது மறுவரிசைப்படுத்திக் கையாளும்;
அல்லது விதியின் முடிவில் ஒரு கூட்டு !!{formula}வை அறிமுகப்படுத்தும்.
எடுத்துக்காட்டாக, $!A, \Gamma \Sequent \Delta$ இலிருந்து $!A \land !B,
\Gamma \Sequent \Delta$ ஐ அனுமானிக்க $\LeftR{\land}$ விதி அனுமதிக்கிறது;
எந்த $\Gamma$, $\Delta$, $!A$, மற்றும்~$!B$ க்கும், $!A, \Gamma \Sequent
\Delta, !B$ இலிருந்து $\Gamma \Sequent \Delta, !A \lif !B$ ஐ அனுமானிக்க
$\RightR{\lif}$ அனுமதிக்கிறது. (குறிப்பாக, $\Gamma$ மற்றும்~$\Delta$
வெறுமையாக இருக்கலாம்.)

% SEGMENT OLP-0065-S02
தொடரணிக் கணியத்தை அடிப்படையாகக் கொண்ட $\Proves$ உறவு பின்வருமாறு
வரையறுக்கப்படுகிறது: $\Gamma_0$ இல் உள்ள ஒவ்வொரு $!A$ உம்~$\Gamma$ இல்
இருக்குமாறும், $\Gamma_0 \Sequent !A$ என்ற தொடரணியை வேராகக் கொண்ட
!!a{derivation} இருக்குமாறும் ஏதோ ஒரு தொடர் $\Gamma_0$ இருந்தால், அப்போதும்
அப்போது மட்டுமே, $\Gamma \Proves !A$. தொடரணிக் கணியத்தில் $!A$ ஒரு
தேற்றம் என்பது, அப்போதும் அப்போது மட்டுமே, $\Sequent !A$ என்ற தொடரணிக்கு
!!a{derivation} உண்டு. எடுத்துக்காட்டாக, $\Proves (!A \land !B) \lif !A$
என்பதைக் காட்டும் !!a{derivation} இதோ:
\begin{prooftree}
  \Axiom$!A \fCenter !A$
  \RightLabel{\LeftR{\land}}
  \UnaryInf$!A \land !B \fCenter !A$
  \RightLabel{\RightR{\lif}}
  \UnaryInf$\fCenter (!A \land !B) \lif !A$
\end{prooftree}

% SEGMENT OLP-0065-S03
ஒவ்வொரு $!A \in \Gamma_0$ உம்~$\Gamma$ இல் உள்ளதாகவும், தொடரணியின்
வலப்புறம் வெறுமையாக உள்ள $\Gamma_0 \Sequent$ க்கு !!a{derivation}
உள்ளதாகவும் இருந்தால்,
கணம் $\Gamma$ தொடரணிக் கணியத்தில் முரணுடையது. \RightR{\Weakening}
விதியைப் பயன்படுத்தி, முரணுடைய கணத்திலிருந்து எந்த !!{sentence}வையும்
!!{derive} இயலும்.

% SEGMENT OLP-0065-S04
தொடரணிக் கணியத்தை 1930களில் கெர்ஹார்ட் கென்ட்சன் உருவாக்கினார். அதன்
முறையான, சமச்சீரான வடிவமைப்பினால், !!{derivation}s பற்றிய கோட்பாட்டை
உருவாக்க அது மிகவும் பயனுள்ள வடிவமுறை. தொடரணிக் கணியத்தில்
!!{derivation}s ஐக் கண்டறிவது ஒப்பீட்டளவில் எளிது; ஆனால்
!!{derivation}s ஐ வாசிப்பது பெரும்பாலும் கடினம், மேலும் நிறுவல்களுடன்
அவற்றுக்குள்ள தொடர்பையும்
சிலவேளைகளில் எளிதாகக் காண முடியாது. இருப்பினும், அது !!{derivation}
முறைமைகளுக்கான மிகவும் நேர்த்தியான அணுகுமுறையாகத் திகழ்ந்துள்ளது; பல
தருக்கவியல்களுக்கு தொடரணிக் கணிய முறைமைகள் உள்ளன.

\end{document}
